\documentclass{article}
\usepackage[T1]{fontenc}
\usepackage{lmodern}
\usepackage{graphicx}
\usepackage{amsmath,amssymb}
\usepackage[a4paper,margin=2cm]{geometry}
\usepackage[absolute,overlay]{textpos}
\setlength{\TPHorizModule}{1mm}
\setlength{\TPVertModule}{1mm}
\usepackage[indonesian]{babel}
\usepackage{hyperref}
\setlength{\emergencystretch}{2em}
% unit: Halaman 103

\begin{document}

% === Halaman 103 (python/native) ===

103

• Dekripsi:


P = K-1 C

 Cipherteks: LNS  atau C = (11, 13, 18) 

 Plainteks:


C = (15, 0, 24) = (P, A, Y) 





















=





















=









































24

0

15

26

mod

570

494

431

18

13

11

17

0

24

6

17

15

15

9

4

\end{document}
